\documentclass[12pt]{ximera}
\typeout{Start loading xmPreamble.tex}%

% Add here extra macro's that are loaded automatically by all documents of claas 'ximera' or 'xourse' in this repo

%%
%%  Example:
%%
% \newcommand{\R}{\mathbb{R}}
\usepackage{tikz}
\usetikzlibrary{positioning, arrows.meta}
\usepackage{multicol}
\usepackage{enumitem}
\usepackage{caption}
\usepackage{amsmath}
\usepackage{amsthm}
\usepackage{fontenc}

%% ---------------------------------------------------------------
%% Shared worksheet macros.
%%
%% These came from the Summer 2026 worksheet set, where each worksheet carried its
%% own copy. They have to live here instead: a xourse inlines every activity into a
%% single document, so duplicate \newcommand definitions across activities collide,
%% and the xourse gobbles \input so a per-topic macro file never loads.
%%
%% They use \providecommand, NOT \newcommand. ximera.cls loads this file automatically,
%% and every activity also carries an explicit \typeout{Start loading xmPreamble.tex}%

% Add here extra macro's that are loaded automatically by all documents of claas 'ximera' or 'xourse' in this repo

%%
%%  Example:
%%
% \newcommand{\R}{\mathbb{R}}
\usepackage{tikz}
\usetikzlibrary{positioning, arrows.meta}
\usepackage{multicol}
\usepackage{enumitem}
\usepackage{caption}
\usepackage{amsmath}
\usepackage{amsthm}
\usepackage{fontenc}

%% ---------------------------------------------------------------
%% Shared worksheet macros.
%%
%% These came from the Summer 2026 worksheet set, where each worksheet carried its
%% own copy. They have to live here instead: a xourse inlines every activity into a
%% single document, so duplicate \newcommand definitions across activities collide,
%% and the xourse gobbles \input so a per-topic macro file never loads.
%%
%% They use \providecommand, NOT \newcommand. ximera.cls loads this file automatically,
%% and every activity also carries an explicit \typeout{Start loading xmPreamble.tex}%

% Add here extra macro's that are loaded automatically by all documents of claas 'ximera' or 'xourse' in this repo

%%
%%  Example:
%%
% \newcommand{\R}{\mathbb{R}}
\usepackage{tikz}
\usetikzlibrary{positioning, arrows.meta}
\usepackage{multicol}
\usepackage{enumitem}
\usepackage{caption}
\usepackage{amsmath}
\usepackage{amsthm}
\usepackage{fontenc}

%% ---------------------------------------------------------------
%% Shared worksheet macros.
%%
%% These came from the Summer 2026 worksheet set, where each worksheet carried its
%% own copy. They have to live here instead: a xourse inlines every activity into a
%% single document, so duplicate \newcommand definitions across activities collide,
%% and the xourse gobbles \input so a per-topic macro file never loads.
%%
%% They use \providecommand, NOT \newcommand. ximera.cls loads this file automatically,
%% and every activity also carries an explicit \input{../xmPreamble}, so this file is
%% read twice. That was harmless while it held only \usepackage lines (LaTeX skips a
%% package it has already loaded) but a second \newcommand is a hard error.
%%
%%   \blank[width]        fill-in-the-blank rule (default 2.5cm)
%%   \showwork            "show and explain your work" reminder
%%   \showworkline        ... plus which schedule line was used
%%   \showworkyear        ... plus which year's schedule was used
%%   \schedhead{...}      centred bold caption above a tiered-rate table
%%   \samesquares[scale]{left}{right}   two equal-area squares, labelled
%%   \samecubes[scale]{left}{right}     two equal-volume cubes, labelled
%% ---------------------------------------------------------------

\providecommand{\blank}[1][2.5cm]{\underline{\hspace{#1}}}
\providecommand{\showwork}{\textbf{REMEMBER TO SHOW AND EXPLAIN YOUR WORK.}}
\providecommand{\showworkline}{\textbf{REMEMBER TO SHOW AND EXPLAIN YOUR WORK.
  Explanation should include how and why you know which line of the
  schedule was used to calculate this amount.}}
\providecommand{\showworkyear}{\begin{sloppypar}\textbf{REMEMBER TO SHOW AND
  EXPLAIN YOUR WORK. Explanation should include how and why you know which
  line of the year's tax schedule was used to calculate this tax
  amount.}\end{sloppypar}}
\providecommand{\schedhead}[1]{\begin{center}\textbf{#1}\end{center}}

\providecommand{\samesquares}[3][1]{%
\begin{center}
{\footnotesize These squares are the \textbf{same size}.}

\vspace{1.5mm}
\begin{tikzpicture}[scale=#1,every node/.style={scale=#1}]
  \draw (0,0) rectangle (2.4,2.4);
  \node[anchor=north west] at (0.15,2.25) {Area:};
  \node[anchor=east]  at (-0.15,1.2) {#2};
  \node[anchor=north] at (1.2,-0.15) {#2};
  \begin{scope}[xshift=4.6cm]
    \draw (0,0) rectangle (2.4,2.4);
    \node[anchor=north west] at (0.15,2.25) {Area:};
    \node[anchor=east]  at (-0.15,1.2) {#3};
    \node[anchor=north] at (1.2,-0.15) {#3};
  \end{scope}
\end{tikzpicture}
\end{center}}

\providecommand{\samecubes}[3][1]{%
\begin{center}
{\footnotesize These cubes are the \textbf{same size}.}

\vspace{1.5mm}
\begin{tikzpicture}[scale=#1,every node/.style={scale=#1}]
  \foreach \dx in {0,5.2} {
    \begin{scope}[xshift=\dx cm]
      \draw (0,0) rectangle (2.0,2.0);
      \draw (0,2.0) -- (0.65,2.6) -- (2.65,2.6) -- (2.0,2.0);
      \fill[black!12] (2.0,0) -- (2.65,0.6) -- (2.65,2.6) -- (2.0,2.0) -- cycle;
      \draw (2.0,0) -- (2.65,0.6) -- (2.65,2.6) -- (2.0,2.0);
      \node[anchor=north west] at (0.1,1.9) {Volume:};
    \end{scope}
  }
  \node[anchor=east]  at (-0.15,1.0) {#2};
  \node[anchor=north] at (1.0,-0.15) {#2};
  \node[anchor=west]  at (2.25,0.4) {#2};
  \node[anchor=east]  at (5.05,1.0) {#3};
  \node[anchor=north] at (6.2,-0.15) {#3};
  \node[anchor=west]  at (7.45,0.4) {#3};
\end{tikzpicture}
\end{center}}

%% NOTE: do NOT add \usepackage to individual activity .tex files.
%% When a xourse inlines an activity it gobbles \documentclass and \input, but a bare
%% \usepackage still executes -- after \begin{document} -- and errors out.
%% All shared packages and macros belong here.
%%
%% ximera.cls already loads: amsmath, amssymb, amsthm, cancel, enumitem, graphicx,
%% hyperref, listings, pgfplots, tikz, url, xcolor. No need to re-load those.
, so this file is
%% read twice. That was harmless while it held only \usepackage lines (LaTeX skips a
%% package it has already loaded) but a second \newcommand is a hard error.
%%
%%   \blank[width]        fill-in-the-blank rule (default 2.5cm)
%%   \showwork            "show and explain your work" reminder
%%   \showworkline        ... plus which schedule line was used
%%   \showworkyear        ... plus which year's schedule was used
%%   \schedhead{...}      centred bold caption above a tiered-rate table
%%   \samesquares[scale]{left}{right}   two equal-area squares, labelled
%%   \samecubes[scale]{left}{right}     two equal-volume cubes, labelled
%% ---------------------------------------------------------------

\providecommand{\blank}[1][2.5cm]{\underline{\hspace{#1}}}
\providecommand{\showwork}{\textbf{REMEMBER TO SHOW AND EXPLAIN YOUR WORK.}}
\providecommand{\showworkline}{\textbf{REMEMBER TO SHOW AND EXPLAIN YOUR WORK.
  Explanation should include how and why you know which line of the
  schedule was used to calculate this amount.}}
\providecommand{\showworkyear}{\begin{sloppypar}\textbf{REMEMBER TO SHOW AND
  EXPLAIN YOUR WORK. Explanation should include how and why you know which
  line of the year's tax schedule was used to calculate this tax
  amount.}\end{sloppypar}}
\providecommand{\schedhead}[1]{\begin{center}\textbf{#1}\end{center}}

\providecommand{\samesquares}[3][1]{%
\begin{center}
{\footnotesize These squares are the \textbf{same size}.}

\vspace{1.5mm}
\begin{tikzpicture}[scale=#1,every node/.style={scale=#1}]
  \draw (0,0) rectangle (2.4,2.4);
  \node[anchor=north west] at (0.15,2.25) {Area:};
  \node[anchor=east]  at (-0.15,1.2) {#2};
  \node[anchor=north] at (1.2,-0.15) {#2};
  \begin{scope}[xshift=4.6cm]
    \draw (0,0) rectangle (2.4,2.4);
    \node[anchor=north west] at (0.15,2.25) {Area:};
    \node[anchor=east]  at (-0.15,1.2) {#3};
    \node[anchor=north] at (1.2,-0.15) {#3};
  \end{scope}
\end{tikzpicture}
\end{center}}

\providecommand{\samecubes}[3][1]{%
\begin{center}
{\footnotesize These cubes are the \textbf{same size}.}

\vspace{1.5mm}
\begin{tikzpicture}[scale=#1,every node/.style={scale=#1}]
  \foreach \dx in {0,5.2} {
    \begin{scope}[xshift=\dx cm]
      \draw (0,0) rectangle (2.0,2.0);
      \draw (0,2.0) -- (0.65,2.6) -- (2.65,2.6) -- (2.0,2.0);
      \fill[black!12] (2.0,0) -- (2.65,0.6) -- (2.65,2.6) -- (2.0,2.0) -- cycle;
      \draw (2.0,0) -- (2.65,0.6) -- (2.65,2.6) -- (2.0,2.0);
      \node[anchor=north west] at (0.1,1.9) {Volume:};
    \end{scope}
  }
  \node[anchor=east]  at (-0.15,1.0) {#2};
  \node[anchor=north] at (1.0,-0.15) {#2};
  \node[anchor=west]  at (2.25,0.4) {#2};
  \node[anchor=east]  at (5.05,1.0) {#3};
  \node[anchor=north] at (6.2,-0.15) {#3};
  \node[anchor=west]  at (7.45,0.4) {#3};
\end{tikzpicture}
\end{center}}

%% NOTE: do NOT add \usepackage to individual activity .tex files.
%% When a xourse inlines an activity it gobbles \documentclass and \input, but a bare
%% \usepackage still executes -- after \begin{document} -- and errors out.
%% All shared packages and macros belong here.
%%
%% ximera.cls already loads: amsmath, amssymb, amsthm, cancel, enumitem, graphicx,
%% hyperref, listings, pgfplots, tikz, url, xcolor. No need to re-load those.
, so this file is
%% read twice. That was harmless while it held only \usepackage lines (LaTeX skips a
%% package it has already loaded) but a second \newcommand is a hard error.
%%
%%   \blank[width]        fill-in-the-blank rule (default 2.5cm)
%%   \showwork            "show and explain your work" reminder
%%   \showworkline        ... plus which schedule line was used
%%   \showworkyear        ... plus which year's schedule was used
%%   \schedhead{...}      centred bold caption above a tiered-rate table
%%   \samesquares[scale]{left}{right}   two equal-area squares, labelled
%%   \samecubes[scale]{left}{right}     two equal-volume cubes, labelled
%% ---------------------------------------------------------------

\providecommand{\blank}[1][2.5cm]{\underline{\hspace{#1}}}
\providecommand{\showwork}{\textbf{REMEMBER TO SHOW AND EXPLAIN YOUR WORK.}}
\providecommand{\showworkline}{\textbf{REMEMBER TO SHOW AND EXPLAIN YOUR WORK.
  Explanation should include how and why you know which line of the
  schedule was used to calculate this amount.}}
\providecommand{\showworkyear}{\begin{sloppypar}\textbf{REMEMBER TO SHOW AND
  EXPLAIN YOUR WORK. Explanation should include how and why you know which
  line of the year's tax schedule was used to calculate this tax
  amount.}\end{sloppypar}}
\providecommand{\schedhead}[1]{\begin{center}\textbf{#1}\end{center}}

\providecommand{\samesquares}[3][1]{%
\begin{center}
{\footnotesize These squares are the \textbf{same size}.}

\vspace{1.5mm}
\begin{tikzpicture}[scale=#1,every node/.style={scale=#1}]
  \draw (0,0) rectangle (2.4,2.4);
  \node[anchor=north west] at (0.15,2.25) {Area:};
  \node[anchor=east]  at (-0.15,1.2) {#2};
  \node[anchor=north] at (1.2,-0.15) {#2};
  \begin{scope}[xshift=4.6cm]
    \draw (0,0) rectangle (2.4,2.4);
    \node[anchor=north west] at (0.15,2.25) {Area:};
    \node[anchor=east]  at (-0.15,1.2) {#3};
    \node[anchor=north] at (1.2,-0.15) {#3};
  \end{scope}
\end{tikzpicture}
\end{center}}

\providecommand{\samecubes}[3][1]{%
\begin{center}
{\footnotesize These cubes are the \textbf{same size}.}

\vspace{1.5mm}
\begin{tikzpicture}[scale=#1,every node/.style={scale=#1}]
  \foreach \dx in {0,5.2} {
    \begin{scope}[xshift=\dx cm]
      \draw (0,0) rectangle (2.0,2.0);
      \draw (0,2.0) -- (0.65,2.6) -- (2.65,2.6) -- (2.0,2.0);
      \fill[black!12] (2.0,0) -- (2.65,0.6) -- (2.65,2.6) -- (2.0,2.0) -- cycle;
      \draw (2.0,0) -- (2.65,0.6) -- (2.65,2.6) -- (2.0,2.0);
      \node[anchor=north west] at (0.1,1.9) {Volume:};
    \end{scope}
  }
  \node[anchor=east]  at (-0.15,1.0) {#2};
  \node[anchor=north] at (1.0,-0.15) {#2};
  \node[anchor=west]  at (2.25,0.4) {#2};
  \node[anchor=east]  at (5.05,1.0) {#3};
  \node[anchor=north] at (6.2,-0.15) {#3};
  \node[anchor=west]  at (7.45,0.4) {#3};
\end{tikzpicture}
\end{center}}

%% NOTE: do NOT add \usepackage to individual activity .tex files.
%% When a xourse inlines an activity it gobbles \documentclass and \input, but a bare
%% \usepackage still executes -- after \begin{document} -- and errors out.
%% All shared packages and macros belong here.
%%
%% ximera.cls already loads: amsmath, amssymb, amsthm, cancel, enumitem, graphicx,
%% hyperref, listings, pgfplots, tikz, url, xcolor. No need to re-load those.

\title{Federal Income Taxes}
\begin{document}
\begin{abstract}
What a tax bracket actually is. Reconstructing the missing base amounts in the 2019
Single tax schedule, computing effective tax rates two different ways, working out the
marginal rate on a raise (and whether a raise can ever leave you worse off), and then
reversing the whole calculation to recover a taxable income from a tax bill.
\end{abstract}
\maketitle

\section*{Setting the Scene}

You may have heard or read news stories about ``tax brackets'' of $12\%$, $24\%$, or
$37\%$. In this activity, we investigate what these percentages mean and how they relate
to federal income taxes.

This activity considers only federal income tax, ignoring other taxes such as Social
Security tax, Medicare tax, and state and local taxes. To keep things simple, we
restrict our attention to \textbf{an unmarried individual whose income sources are
entirely wages, salaries, and tips, and who is not claimed as a dependent on anyone
else's tax return.}

\textbf{This is not a tax class, and your instructor is not a tax specialist. If you
have any questions or concerns about your taxes or filing your taxes, you should talk
to a tax professional.} We are looking only at the mathematics behind the calculation.

For this example, let's assume a fictional person has an income of \$75,000. The 1040
instructions used here can be found at
\url{https://www.irs.gov/pub/irs-prior/i1040gi--2019.pdf}.

Where exactly do the values in the Tax Table come from? The entire tax table can be
described more concisely by something called a \textbf{Tax Schedule}: a table of
instructions for how to calculate taxes based on the taxable income value. The Tax Table
values are calculated from these instructions for the taxpayers' convenience. Notice
that for anyone with a taxable income over \$100,000, the table does not have taxes
pre-calculated, so those individuals must follow the instructions from the Tax Schedule.

\section*{The 2019 Tax Schedule}

\schedhead{Single Taxable Income Tax Brackets and Rates, 2019}

\begin{center}
\begin{tabular}{|l|l|l|}
\hline
\textbf{Over} & \textbf{But not over} & \textbf{The tax is} \\
\hline
\$0        & \$9,700    & $10\%$ of the taxable income \\
\hline
\$9,700    & \$39,475   & \textbf{A} plus $12\%$ of the amount over \$9,700 \\
\hline
\$39,475   & \$84,200   & \textbf{B} plus $22\%$ of the amount over \$39,475 \\
\hline
\$84,200   & \$160,725  & \textbf{C} plus $24\%$ of the amount over \$84,200 \\
\hline
\$160,725  & \$204,100  & \$32,748 plus $32\%$ of the amount over \$160,725 \\
\hline
\$204,100  & \$510,300  & \$46,628 plus $35\%$ of the amount over \$204,100 \\
\hline
\$510,300  & no limit   & \$153,798 plus $37\%$ of the amount over \$510,300 \\
\hline
\end{tabular}
\end{center}

Notice that there are blanks in the Tax Schedule. We are going to fill in these values
as a way of helping us understand the instructions.

\subsection*{Blank A}

\begin{problem}
If someone had a taxable income of \$9,700, how much would they pay in taxes according
to the schedule?

\$$\answer{970}$

\begin{solution}
\$9,700 falls in the first line of the schedule, so the tax is $10\%$ of the taxable
income:
\[
0.10 \cdot \$9{,}700 = \$970.
\]
\end{solution}
\end{problem}

\begin{problem}
If someone had a taxable income of \$9,800, how much would they pay in taxes according
to the schedule?

\$$\answer{982}$

\begin{freeResponse}
\end{freeResponse}

\begin{solution}
\$9,800 falls in the second line, so the tax is A plus $12\%$ of the amount over
\$9,700. The amount over \$9,700 is \$100, so the tax is
\[
\text{A} + 0.12 \cdot \$100 = \text{A} + \$12.
\]
We do not yet know A --- but we do know from the previous problem that someone at
exactly \$9,700 pays \$970.
\end{solution}
\end{problem}

\begin{problem}
What can we put in the blank space for A, and why?

A $=$ \$$\answer{970}$

\begin{freeResponse}
\end{freeResponse}

\begin{solution}
A $=$ \$970.

The reasoning is continuity. At a taxable income of exactly \$9,700 the first line gives
\$970. The second line applies to income \emph{over} \$9,700, and at \$9,700 the ``amount
over \$9,700'' is \$0 --- so the second line must also give \$970 there, meaning
A $+ 0.12 \cdot \$0 = \$970$, so A $= \$970$.

Put plainly: A is simply all the tax already owed on the first \$9,700 of income. The
schedule charges you the full tax on each earlier bracket, then adds the new rate on
only the part of your income that spills into the new bracket.
\end{solution}
\end{problem}

\subsection*{Blank B}

\begin{problem}
If someone had a taxable income of \$39,475, how much would they pay in taxes according
to the schedule?

\$$\answer{4543}$

\begin{solution}
\$39,475 is the top of the second line, so use that line: \$970 plus $12\%$ of the
amount over \$9,700.
\[
\$970 + 0.12 \cdot (\$39{,}475 - \$9{,}700) = \$970 + 0.12 \cdot \$29{,}775
= \$970 + \$3{,}573 = \$4{,}543.
\]
\end{solution}
\end{problem}

\begin{problem}
If someone had a taxable income of \$39,775, how much would they pay in taxes according
to the schedule?

\$$\answer{4609}$

\begin{solution}
\$39,775 falls in the third line, so the tax is B plus $22\%$ of the amount over
\$39,475. That amount over is \$300, and by the same continuity argument as before,
B $=$ \$4,543. So
\[
\$4{,}543 + 0.22 \cdot \$300 = \$4{,}543 + \$66 = \$4{,}609.
\]
\end{solution}
\end{problem}

\begin{problem}
What can we put in the blank space for B, and why?

B $=$ \$$\answer{4543}$

\begin{freeResponse}
\end{freeResponse}

\begin{solution}
B $=$ \$4,543 --- exactly the tax owed by someone at the very top of the previous
bracket, for the same reason as A. It is the total tax accumulated on the first
\$39,475 of taxable income.
\end{solution}
\end{problem}

\subsection*{Blank C}

\begin{problem}
Describe the pattern you notice from the previous two problems, then calculate the value
for blank C.

C $=$ \$$\answer[tolerance=0.5]{14382.50}$

\begin{freeResponse}
\end{freeResponse}

\begin{solution}
The pattern: each blank is the total tax owed by someone sitting exactly at the bottom
edge of that bracket --- equivalently, the tax accumulated across all the brackets
below it.

So C is the tax on a taxable income of \$84,200, computed from the third line:
\[
\text{C} = \$4{,}543 + 0.22 \cdot (\$84{,}200 - \$39{,}475)
= \$4{,}543 + 0.22 \cdot \$44{,}725 = \$4{,}543 + \$9{,}839.50 = \$14{,}382.50.
\]

You can check the pattern against a line the schedule already fills in. The next base
should be
\[
\$14{,}382.50 + 0.24 \cdot (\$160{,}725 - \$84{,}200) = \$14{,}382.50 + \$18{,}366
= \$32{,}748.50,
\]
and the schedule prints \$32,748 --- matching, with the IRS rounding to whole dollars.
\end{solution}
\end{problem}

\section*{Effective Tax Rate}

There are two very important percentages related to federal income taxes. The first is
the \textbf{effective tax rate}.

\begin{problem}
Return to our original example of the person with \$75,000 of total income, taking the
2019 standard deduction of \$12,200. Calculate their tax from the Tax Schedule.

Taxable income: \$$\answer{62800}$

Tax: \$$\answer[tolerance=1]{9674.50}$

\begin{solution}
Taxable income is the AGI minus the standard deduction:
\[
\$75{,}000 - \$12{,}200 = \$62{,}800.
\]
That falls in the third line of the schedule, so
\[
\$4{,}543 + 0.22 \cdot (\$62{,}800 - \$39{,}475) = \$4{,}543 + 0.22 \cdot \$23{,}325
= \$4{,}543 + \$5{,}131.50 = \$9{,}674.50.
\]
\end{solution}
\end{problem}

\begin{problem}
What percent of this person's \textbf{TOTAL} income went to taxes?

About $\answer[tolerance=0.3]{12.9}$\%

\begin{solution}
\[
\frac{\$9{,}674.50}{\$75{,}000} = 0.12899\ldots \approx 12.9\%.
\]
\end{solution}
\end{problem}

\begin{problem}
What percent of this person's \textbf{TAXABLE} income went to taxes?

About $\answer[tolerance=0.3]{15.4}$\%

\begin{solution}
\[
\frac{\$9{,}674.50}{\$62{,}800} = 0.15405\ldots \approx 15.4\%.
\]
The same tax bill, divided by a smaller base, gives a noticeably larger percentage.
\end{solution}
\end{problem}

\begin{problem}
Of the two percentages you just calculated, which do you think better describes how much
this person was taxed? And can you think of a reason why someone might want the other
one?

\begin{freeResponse}
\end{freeResponse}

\begin{solution}
Either can be defended, which is the point.

The percentage of \emph{total} income ($12.9\%$) is arguably the better description of
what actually happened to this person: they earned \$75,000 and \$9,674.50 of it went to
federal income tax. That is the figure that answers ``how much of what I made did I hand
over?''

The percentage of \emph{taxable} income ($15.4\%$) is the better description of how the
tax system itself behaved, because taxable income is the number the schedule actually
operates on. Someone comparing the burden of the schedule across people --- or checking
the schedule's own arithmetic --- would want this one, since the standard deduction is a
separate policy choice layered on top.

BOTH values can be referred to as an \textbf{effective tax rate}. Different
professionals use different comparisons depending on what they are investigating or
communicating. They are both telling you, in some sense, effectively how much of a
person's income is being taxed as a whole. When you read an effective tax rate quoted
anywhere, it is worth asking which denominator was used.
\end{solution}
\end{problem}

\section*{Marginal Tax Rates}

The second important percentage is the \textbf{marginal tax rate}. This one is more
difficult to wrap your head around, and is the one that causes perhaps the most common
confusion related to taxes: \textbf{will entering a new tax bracket cost me more money?}

A \textbf{tax bracket} is each window of income listed in the Tax Schedule. Taxable
income from \$0 to \$9,700 is the first tax bracket; \$9,700 to \$39,475 is the next;
and so on. There is a common concern that if someone with a taxable income of \$9,000
gets a raise to \$10,000, they will end up taking home \emph{less} money after taxes
than before the raise.

A \textbf{marginal tax rate} is what percentage of a \emph{portion} of income is paid as
taxes. This differs from the effective tax rate, which looks at income as a whole.

\begin{problem}
Return to your work for Blank A. Tax on \$9,700 of taxable income was \$970; tax on
\$9,800 was \$982. What was the marginal tax rate on the last \$100 of income --- that
is, what percentage of that last \$100 was paid in taxes?

$\answer{12}$\%

\begin{freeResponse}
\end{freeResponse}

\begin{solution}
The extra tax is $\$982 - \$970 = \$12$, paid on \$100 of extra income:
\[
\frac{\$12}{\$100} = 0.12 = 12\%.
\]

This is exactly the $12\%$ printed in the second line of the schedule, and that is not a
coincidence. The schedule's rate for a bracket \emph{is} the marginal rate on income
falling inside that bracket --- that is what the instruction ``plus $12\%$ of the amount
over \$9,700'' is saying.
\end{solution}
\end{problem}

\begin{problem}
Now let's see how this relates to raises. Someone has a taxable income of \$37,000 and
is offered a raise that would take them to \$42,000 --- crossing a bracket boundary.

Their taxable income increases by \$$\answer{5000}$.

Tax at \$37,000: \$$\answer[tolerance=1]{4246}$

Tax at \$42,000: \$$\answer[tolerance=1]{5098.50}$

They would pay \$$\answer[tolerance=1]{852.50}$ more in taxes.

\begin{solution}
The raise is $\$42{,}000 - \$37{,}000 = \$5{,}000$.

At \$37,000 (second line):
\[
\$970 + 0.12 \cdot (\$37{,}000 - \$9{,}700) = \$970 + \$3{,}276 = \$4{,}246.
\]

At \$42,000 (third line):
\[
\$4{,}543 + 0.22 \cdot (\$42{,}000 - \$39{,}475) = \$4{,}543 + \$555.50 = \$5{,}098.50.
\]

The extra tax is $\$5{,}098.50 - \$4{,}246 = \$852.50$.
\end{solution}
\end{problem}

\begin{problem}
What is the marginal tax rate on their raise --- what percentage of the raise goes to
taxes? How is this related to the Tax Schedule instructions, and should they take the
raise?

About $\answer[tolerance=0.3]{17.05}$\%

\begin{freeResponse}
\end{freeResponse}

\begin{solution}
\[
\frac{\$852.50}{\$5{,}000} = 0.1705 = 17.05\%.
\]

Notice this is neither $12\%$ nor $22\%$, but sits between them. That is because the
raise straddles a bracket boundary: the first \$2,475 of it (from \$37,000 up to
\$39,475) is taxed at $12\%$, and the remaining \$2,525 (from \$39,475 to \$42,000) is
taxed at $22\%$. The blended result is $17.05\%$.

\textbf{Should they take the raise? Yes, unambiguously.} They pay \$852.50 more in tax
on \$5,000 more income, so they keep $\$5{,}000 - \$852.50 = \$4{,}147.50$ more than
before. Take-home pay went \emph{up}.

This is the answer to the common fear. Entering a higher bracket never reduces your
take-home pay, because the higher rate applies only to the portion of income above the
bracket boundary --- not retroactively to all of it. For a raise to leave you worse off,
a marginal rate would have to exceed $100\%$, and no line of the schedule comes close.
\end{solution}
\end{problem}

\section*{Working ``Backwards''}

People like to guard their tax filing information, not just for privacy reasons, but
because all of the calculations we do from the schedule are \textbf{reversible}. If we
know how much someone paid in taxes, we can determine what they claimed as taxable
income. We continue with the same assumptions: a single person, standard deduction,
2019.

\begin{problem}
Suppose someone paid \$6,700 in taxes. What line of the Tax Schedule must have been used
to calculate this tax amount? Explain how you know it must be this line by explaining
why it could not have been the previous line nor the next line.

\begin{freeResponse}
\end{freeResponse}

\begin{solution}
It must be the third line, for taxable income over \$39,475 but not over \$84,200.

It cannot be the second line: the most tax that line can produce is at its top edge,
\$39,475, which gives \$4,543. Since $\$6{,}700 > \$4{,}543$, the second line cannot
reach this tax amount.

It cannot be the fourth line either: that line starts at \$14,382.50 (blank C, the tax
at \$84,200). Since $\$6{,}700 < \$14{,}382.50$, the fourth line cannot produce a tax
this small.

So \$6,700 sits between \$4,543 and \$14,382.50, which brackets it squarely in the third
line. This is the general method: each line covers a known \emph{range of tax amounts},
and you locate which range the tax bill falls into.
\end{solution}
\end{problem}

\begin{problem}
Reverse the steps of that line to determine what the taxable income must have been.

About \$$\answer[tolerance=50]{49279.55}$

\begin{solution}
The third line says $\text{tax} = \$4{,}543 + 0.22 \cdot (T - \$39{,}475)$. Setting that
equal to \$6,700 and undoing each step:
\[
\$6{,}700 - \$4{,}543 = \$2{,}157 = 0.22 \cdot (T - \$39{,}475)
\]
\[
T - \$39{,}475 = \frac{\$2{,}157}{0.22} = \$9{,}804.55
\]
\[
T = \$39{,}475 + \$9{,}804.55 = \$49{,}279.55.
\]
\end{solution}
\end{problem}

\begin{problem}
Write a function for this line of the tax schedule --- that is, for a taxable income
$T$, an expression $\text{tax}(T)$ that calculates the tax. How could you use that
function to answer the previous question? And if this person took the \$12,200 standard
deduction, what was their total income?

Total income: about \$$\answer[tolerance=50]{61479.55}$

\begin{freeResponse}
\end{freeResponse}

\begin{solution}
For the third line,
\[
\text{tax}(T) = 4543 + 0.22(T - 39475),
\]
which simplifies to $\text{tax}(T) = 0.22T - 4141.5$.

Written this way, the earlier question becomes: solve $\text{tax}(T) = 6700$ for $T$.
Reversing the arithmetic step by step, as we did above, is exactly solving that
equation --- but having the function makes it clear that this is a single linear
equation in one unknown, and that it has exactly one solution.

Total income is taxable income plus the deduction that was subtracted to get there:
\[
\$49{,}279.55 + \$12{,}200 = \$61{,}479.55.
\]
\end{solution}
\end{problem}

\begin{problem}
Now suppose someone paid \$17,500 in taxes. What was this person's taxable income? Be
sure to support your answer by stating how you know which line of the Tax Schedule was
used, and explaining what you are doing and why in terms of the information in the
schedule.

About \$$\answer[tolerance=100]{97189.58}$

\begin{freeResponse}
\end{freeResponse}

\begin{solution}
First identify the line. The fourth line runs from a tax of \$14,382.50 (at a taxable
income of \$84,200) up to \$32,748.50 (at \$160,725). Since
$\$14{,}382.50 < \$17{,}500 < \$32{,}748.50$, the fourth line was used.

That line says $\text{tax} = \$14{,}382.50 + 0.24 \cdot (T - \$84{,}200)$. Reversing:
\[
\$17{,}500 - \$14{,}382.50 = \$3{,}117.50 = 0.24 \cdot (T - \$84{,}200)
\]
\[
T - \$84{,}200 = \frac{\$3{,}117.50}{0.24} = \$12{,}989.58
\]
\[
T = \$84{,}200 + \$12{,}989.58 = \$97{,}189.58.
\]

Each step undoes one instruction from the schedule: subtracting the base amount removes
the tax already owed on the lower brackets, and dividing by the rate converts the
remaining tax back into the income it was charged on.
\end{solution}
\end{problem}

\end{document}
